\subsection{线程束停顿与延迟隐藏}

线程束状态统计（Warp State Statistics）把平均线程束周期分配到被选中、就绪但未被选中以及各类停顿状态。\term{记分板}{scoreboard}跟踪指令依赖是否已经满足：如果下一条指令需要某个尚未返回的加载结果，该线程束就必须等待。其他线程束仍可继续执行，因此某类停顿存在，并不自动说明整个 SM 同时空闲；应先结合调度报告确认是否有大量未使用的发射机会。

表\ref{l14:tab:stalls}列出与矩阵乘法诊断最相关的几类状态。对应硬件通路的名称较具体，尤其需要区分等待数据返回与等待队列空位。

\begin{table}[htbp]
\centering\small
\caption{常见线程束等待状态及其解释}
\label{l14:tab:stalls}
\begin{tabularx}{\textwidth}{@{}p{43mm}YY@{}}
\toprule
状态 & 正在等待什么 & 可进一步检查的内容\\
\midrule
长记分板等待\newline（Long Scoreboard） & L1／纹理路径操作产生的数据依赖，包括全局、局部、纹理或表面访问 & 请求合并、缓存命中、加载到使用之间的独立工作\\
短记分板等待\newline（Short Scoreboard） & 非 L1／纹理路径的相关操作依赖；共享读取是常见原因 & 共享布局、存储体冲突及共享数据的消费位置\\
存储输入输出节流\newline（MIO Throttle） & 对应指令队列尚无空位 & 共享加载数量、可否使用较宽加载，以及相关通路压力\\
数学流水线节流\newline（Math Pipe Throttle） & 所需数学流水线暂时不能接收下一条指令 & 指令是否集中在特定流水线，有无冗余运算或可调整的指令组合\\
同步屏障等待\newline（Barrier） & 同步要求的其他线程束尚未到达 & 块内工作是否均衡、加载与计算分工、同步的位置\\
未被选中\newline（Not Selected） & 线程束已经就绪，但本周期选择了其他线程束 & 结合实际发射率判断；存在这种状态说明调度器有可选工作\\
\bottomrule
\end{tabularx}
\end{table}

长记分板等待涉及的加载可能命中缓存，因而不能直接解释成全部时间都在等待显存。短记分板和存储输入输出通路还可能涉及特殊数学或动态分支等操作，具体应以热点指令为准。存储输入输出节流侧重队列接收能力；如果大量窄共享加载不断挤入同一入口，减少加载指令数量就可能有效，即使共享数据的总字节数保持不变。

在前述简单内核的示例中，每条发射指令对应约 $25.18$ 个线程束周期，其中长记分板等待约占 $20.6$ 周期，即约 $81.8\%$。结合接近满额的驻留数量和较低的就绪数量，可以把重点放在加载依赖及其延迟隐藏上。增加更多驻留线程束受到硬件上限限制；提高数据复用、让加载更高效，或在等待结果期间提供其他独立计算，都可能改善当前状态。

同一线程束中的单一累加器也会形成连续依赖：下一次乘加需要前一次乘加的结果。使用多个输出累加器，可以提供多条独立链，让编译器有机会交错安排指令。这与线程束之间的切换共同隐藏延迟，但会增加寄存器占用。诊断时需要观察最终发射率与总时间，不能把某个占用率或某类停顿的单项变化作为唯一目标。

停顿统计还常以每条发射指令或采样次数作为分母。如果优化删除了很多指令，某种剩余等待的比例可能上升；执行时间仍可能下降，因为总工作量和主要队列压力已经减少。比较时应保持分母意识，同时观察绝对指令数量、运行时间和多个相关指标。

\subsection{从源代码行追踪到机器指令}

源码热点视图将设备代码行号与编译后的指令联系起来。虚拟指令集（Parallel Thread Execution，PTX）提供中间表示，目标架构机器指令（SASS）则反映具体 GPU 上使用的加载、存储、乘加和控制指令。源代码中的一行表达式可能展开为地址计算、加载、乘加与循环控制等多条指令；编译器还会进行展开、重排和寄存器分配，所以源码行与机器指令通常不一一对应。

使用 \code{-lineinfo} 编译后，可以按源码行查看执行指令数量、寄存器活跃情况与停顿采样，再沿 PTX 和 SASS 确认实际执行路径。例如，\code{LDG} 表示全局读取，\code{LDS} 表示共享读取，\code{STS} 表示共享写入，\code{FFMA} 表示单精度融合乘加。指令后缀还可能说明访问宽度或其他执行语义。

一个典型依赖链可以写成下面的机器指令结构。第一条从全局内存加载到寄存器 \code{R2}，第二条乘加使用 \code{R2}；地址寄存器和其他源寄存器在此前已经准备好。

\noindent\begin{minipage}{\linewidth}
\begin{lstlisting}[language={},caption={加载结果到消费指令的依赖示意}]
LDG.E R2, [R4];
FFMA  R8, R2, R6, R8;
\end{lstlisting}
\end{minipage}

若 \code{R2} 尚未准备好，停顿可能记录在等待使用它的乘加指令上。因此，热点落在乘加行并不说明乘加流水线本身一定是瓶颈；需要沿寄存器依赖回溯到生产该数据的加载。类似地，一条共享存储也可能因为等待前面的全局加载结果而呈现长记分板等待。

\term{存活寄存器}{live registers}表示在某个程序位置仍保存着后续需要使用的值的寄存器。它帮助理解哪些临时值同时占据状态，与整个内核最终分配的每线程寄存器数量相关，但两者不是逐行相同的数值。查看向量化或线程粗化后的代码时，应同时确认宽读取是否生成、寄存器压力是否增大，以及占用率与局部内存访问是否发生变化。

一条完整的定位过程由这些观察逐步收敛：时间线确定主要耗时内核，吞吐与屋顶线提供性能上界，存储和调度报告定位繁忙通路或等待类型，源码热点与机器指令再确定可以修改的代码。修改之后，预期受影响的计数器应出现对应变化，并由正常计时验证实际收益。
